\begin{thebibliography}{28}
\bibitem{CarreLascouxLeclerc} Christophe Carr\'e, Alain Lascoux, and Bernard Leclerc, Turbo-straightening for decomposition into standard bases, Internat. J. Algebra Comput. 2 (1992), no. 03, 275--290.
\bibitem{Clifford} A. H. Clifford, Matrix representations of completely simple semigroups, Amer. J. Math. 64 (1942), 327--342.
\bibitem{COSSZ} Laura Colmenarejo, Rosa Orellana, Franco Saliola, Anne Schilling, and Mike Zabrocki, An insertion algorithm on multiset partitions with applications to diagram algebras, J. Algebra 557 (2020), 97--128.
\bibitem{FitzGerald.2003} D. G. FitzGerald, A presentation for the monoid of uniform block permutations, Bull. Austral. Math. Soc. 68 (2003), no. 2, 317--324.
\bibitem{GMS09} Olexandr Ganyushkin, Volodymyr Mazorchuk, and Benjamin Steinberg, On the irreducible representations of a finite semigroup, Proc. Amer. Math. Soc. 137 (2009), no. 11, 3585--3592.
\bibitem{HJ} Tom Halverson and Theodore N. Jacobson, Set-partition tableaux and representations of diagram algebras, Algebr. Comb. 3 (2020), no. 2, 509--538.
\bibitem{Har} Nate Harman, Representations of monomial matrices and restriction from $GL_n$ to $S_n$, Apr 2018, \url{http://arxiv.org/abs/1804.04702}.
\bibitem{OEIS} OEIS Foundation Inc., The On-Line Encyclopedia of Integer Sequences, 2019, \url{http://oeis.org}, [Online].
\bibitem{Jones.1994} V. F. R. Jones, The Potts model and the symmetric group, in Subfactors (Kyuzeso, 1993), World Sci. Publ., River Edge, NJ, 1994, pp. 259--267.
\bibitem{Kosuda.2000} Masashi Kosuda, Characterization for the party algebras, Ryukyu Math. J. 13 (2000), 7--22.
\bibitem{Kosuda.2001} Masashi Kosuda, Party algebra and construction of its irreducible representations, Formal Power Series and Algebraic Combinatorics (FPSAC01), Tempe, Arizona (USA), 2001, pp. 20--26.
\bibitem{Kosuda.2006} Masashi Kosuda, Irreducible representations of the party algebra, Osaka J. Math. 43 (2006), no. 2, 431--474.
\bibitem{LP1969} G\'erard Lallement and Mario Petrich, Irreducible matrix representations of finite semigroups, Trans. Amer. Math. Soc. 139 (1969), 393--412.
\bibitem{LR} Nicholas A Loehr and Jeffrey B Remmel, A computational and combinatorial expos\'e of plethystic calculus, J. Algebraic Combin. 33 (2011), no. 2, 163--198.
\bibitem{Martin.2000} P. P. Martin, The partition algebra and the Potts model transfer matrix spectrum in high dimensions, J. Phys. A 33 (2000), no. 19, 3669--3695.
\bibitem{Martin.1991} Paul Martin, Potts models and related problems in statistical mechanics, Series on Advances in Statistical Mechanics, vol. 5, World Scientific Publishing Co., Inc., Teaneck, NJ, 1991.
\bibitem{Martin.1994} Paul Martin, Temperley-Lieb algebras for nonplanar statistical mechanics---the partition algebra construction, J. Knot Theory Ramifications 3 (1994), no. 1, 51--82.
\bibitem{Martin.1996} Paul Martin, The structure of the partition algebras, J. Algebra 183 (1996), no. 2, 319--358.
\bibitem{McAlister} Donald B McAlister, Characters of finite semigroups, J. Algebra 22 (1972), no. 1, 183--200.
\bibitem{Munn} W. D. Munn, On semigroup algebras, Proc. Cambridge Philos. Soc. 51 (1955), 1--15.
\bibitem{Naruse.2005} Hiroshi Naruse, Characters of party algebras, slides of talk for Workshop on Cellular and Diagram Algebras in Mathematics and Physics at Oxford Univ., dated 04-04-2005, 2005, \url{https://www.ccn.yamanashi.ac.jp/~hnaruse/ken/050404OHP.pdf}.
\bibitem{Ponizovski} I. S. Ponizovski\u\i, On matrix representations of associative systems, Mat. Sb. N.S. 38(80) (1956), 241--260.
\bibitem{RZ1991} John Rhodes and Yechezkel Zalcstein, Elementary representation and character theory of finite semigroups and its application, in Monoids and semigroups with applications (Berkeley, CA, 1989), World Sci. Publ., River Edge, NJ, 1991, pp. 334--367.
\bibitem{Sagan.2001} Bruce E. Sagan, The symmetric group: Representations, combinatorial algorithms, and symmetric functions, second ed., Graduate Texts in Mathematics, vol. 203, Springer-Verlag, New York, 2001.
\bibitem{Solomon} Louis Solomon, Representations of the rook monoid, J. Algebra 256 (2002), no. 2, 309--342.
\bibitem{Steinberg-Mobius} Benjamin Steinberg, M\"obius functions and semigroup representation theory II: Character formulas and multiplicities, Adv. Math. 217 (2008), 1521--1557.
\bibitem{Steinberg.2016} Benjamin Steinberg, Representation theory of finite monoids, Universitext, Springer, Cham, 2016.
\bibitem{Tanabe.1997} Kenichiro Tanabe, On the centralizer algebra of the unitary reflection group $G(m,p,n)$, Nagoya Math. J. 148 (1997), 113--126.
\end{thebibliography}
